\slajdid{10-s008}
\begin{frame}[label=10-s008]{Impulsni režim rada bipolarnog tranzistora}
\fontsize{9}{9.5}\selectfont
\begin{columns}[T,onlytextwidth]\column{.255\textwidth}\centering
\begin{tikzpicture}[x=.45cm,y=.32cm,font=\sitneoznake,>=Latex]\ctikzset{bipoles/length=.65cm}\draw(0,-.5)--(0,.5);\draw(-.6,0)--(0,0);\draw(0,.32)--(.7,.7)--(.7,1.3);\draw[->](0,-.32)--(.7,-.7);\draw(.7,-.7)--(.7,-1.3)node[ground]{};\node[right]at(.25,0){T1};\draw(-2.3,0)node[ocirc]{}to[R,l=$R_B$](-.6,0);\draw(.7,1.3)to[R,l_=$R_C$](.7,3);\draw(.4,3)--(1,3);\node[above]at(.7,3){$V_{CC}$};\draw(.7,1.3)--(2,1.3)node[ocirc]{};\draw(-2.3,-1.3)node[ground]{}--(-2.3,-.8);\draw[->](-2.3,-.7)--(-2.3,-.05)node[midway,left]{$v_i(t)$}node[near end,right]{$+$};\draw(2,-.7)node[ground]{}--(2,-.3);\draw[->](2,-.2)--(2,1.25)node[midway,right]{$v_o(t)$}node[near end,right]{$+$};\draw[->](-.65,.65)--(-.15,.65)node[midway,above]{$i_B(t)$};\end{tikzpicture}
\par\smallskip
\begin{tikzpicture}[x=.45cm,y=.3cm,font=\sitneoznake,>=Latex]\draw[->](0,0)--(3.2,0)node[right]{t};\draw[->](.2,-1.1)--(.2,2.1)node[left]{$v_i(t)$};\draw(.2,-.8)--(.7,-.8)--(.7,1.3)--(2.3,1.3)--(2.3,-.8);\node[left]at(.2,1.3){$V_H$};\node[left]at(.2,-.8){$V_L$};\end{tikzpicture}

\par\smallskip
\begin{tikzpicture}[x=.34cm,y=.3cm,font=\sitneoznake,>=Latex]
\draw[->](0,5.5)--(8,5.5)node[below]{t};\draw[->](.2,4.1)--(.2,7.4)node[right]{$v_i(t)$};\draw[dashed](.2,6.7)--(1,6.7);\node[left]at(.2,6.7){$V_H$};\node[left]at(.2,4.3){$V_L$};\draw(.2,4.3)--(1,4.3)--(1,6.7)--(3.5,6.7)--(3.5,4.3)--(7.4,4.3);
\draw[->](0,2.7)--(8,2.7)node[below]{t};\draw[->](.2,2.3)--(.2,4);\node[fill=white,inner sep=1pt]at(3.3,3.7){$v_o(t)$};\node[left]at(.2,3.7){$V_{OH}$};\node[left]at(.2,2.85){$V_{OL}$};\draw[dashed](.2,2.85)--(7.4,2.85);\draw(.2,3.7)--(1.55,3.7)--(2.1,2.85)--(4.5,2.85)--(5.5,3.7)--(7.4,3.7);
\draw[->](0,.4)--(8,.4)node[below]{t};\draw[->](.2,-1.1)--(.2,2)node[right]{$i_B(t)$};\draw(.2,.4)--(1,.4)..controls(1.6,.5)and(1.8,1.1)..(2.4,1.1)--(3.5,1.1)--(3.5,-.45)--(4.5,-.45)..controls(4.65,-.2)and(5.2,-.2)..(5.5,-.2)..controls(5.8,.2)and(6.6,.38)..(7.4,.4);
\foreach\x/\n in{1/0,1.55/1,2.1/2,3.5/3,4.5/4,5.5/5,7.4/6}{\draw[dashed](\x,-.9)--(\x,4);\ifnum\n=1\node[below]at(\x,-1.6){$t_\n$};\else\node[below]at(\x,-.9){$t_\n$};\fi}
\end{tikzpicture}

\column{.725\textwidth}
Period $t_0$ do $t_1$ – vreme kašnjenja; pune se oblasti, „kapacitivnosti“, koje su bile inverzno polarizovane i „skupljaju“ se, i dalje je spoj inverzno polarizovan, još uvek ne postoji tranzistorski efekat, struja kolektora jednaka nuli.
\par
Period $t_1$ do $t_2$ – vreme opadanja; pune se i dalje oblasti prostornog tovara, „kapacitivnosti“, bazno emitorski spoj je direktno polarizovan, uska je oblast, pojavljuje se višak manjinskih nosilaca u bazi, postoji tranzistorski efekat, struja kolektora postoji i zavisi od struje baze. Napon na izlazu opada.
\par
Period $t_2$ do $t_3$ – pri kraju prethodnog perioda i spoj baza emiter i spoj baza kolektor su direktno polarizovani. Pojavio se veliki višak manjinskih nosilaca na oba spoja. Tranzistor je u zasićenju.
\par
Period $t_3$ do $t_4$ – period zasićenja. Napon na ulazu u kolo je negativan. Međutim postoje viškovi manjinskih nosilaca u oblasti baze. Oni ne mogu trenutno de nestanu. Pojavljuje se negativna bazna struja koja eliminiše viškove nosilaca iz baze. Tranzistor je i dalje u zasićenju-
\par
Period $t_4$ do $t_5$ – vreme uspostavljanja. Oblast baza kolektor je postala inverzno polarizovano i nastavlja da se širi. Opada broj nosilaca uz oblast baza emitor, opada kolektorska struja.
\par
Period $t_5$ do $t_6$ – vreme oporavka. Oblast baza kolektor je inverzno polarizovano. Oblast baza emitor postaje inverzno polarizovana i širi se. Pune se oblasti prostornog tovara „kapacitivnosti“. Zato i jeste eksponencijalan oblik.
\end{columns}
\end{frame}
